\section{敏感性分析}

本文区分数值收敛、输入假设和机制因素三类敏感性。Q3 的 $N=512$ 到 $N=1024$ 事件时刻变化为 $0.132\,\mathrm s$，Q4 收缩—附录 4 主情景变化为 $0.136\,\mathrm s$；收紧时间容差后 Q3、Q4 事件变化分别为 $-0.198\,\mathrm s$ 和 $-0.520\,\mathrm s$。

长期计算对环境延拓较敏感：将 4 h 后环境由末小时均值改为最后观测值，Q3 事件提前约 18.28 min，Q4 主情景提前约 16.07 min。Q4 半径轨迹改用 PCHIP 插值时，事件提前约 12.75 s，占总时长约 $7\times10^{-5}$。上述结果说明数值离散误差较小，但长期预测仍受边界假设影响。

当前未对 $h$、$k_m$ 或扩散率系数开展系统的单参数或 Sobol 全局敏感性分析，因此本节不将已完成的收敛和输入对照扩写为完整参数敏感度结论。
